\chapter{Generative UI}
\label{ch:generative-ui}

Generative UI allows agents to render rich, interactive components inline in a chat conversation instead of returning only text. When an agent calls a tool, the tool's result can be rendered as a structured UI component---a table, a card, a diff view, a form---rather than a plain text bubble.

\haira{} approaches generative UI at the language level. Tools declare what component renders their output. The runtime ships a built-in component catalog. No frontend framework, bundler, or client-side JavaScript is required.

\section{Motivation}

Consider a migration tool that returns a dry-run result:

\begin{lstlisting}
tool dry_run(file: string) -> string {
    """Run SQL against the database in dry-run mode"""
    // ... returns "DRY-RUN FAILED on seeds.\n\nROOT CAUSE: ..."
}
\end{lstlisting}

The agent receives a string, streams it as markdown, and the user sees a wall of text. The information is there, but the presentation is poor---error details, affected tables, and fix SQL are all flattened into one block.

With generative UI, the same tool returns structured data and declares a renderer:

\begin{lstlisting}
@render("status-card")
tool dry_run(file: string) -> DryRunResult {
    """Run SQL against the database in dry-run mode"""
    // ... returns structured result
}
\end{lstlisting}

The chat UI renders a status card with colored indicators, collapsible sections, and copy-paste SQL blocks---all from the same tool declaration, with zero frontend code.

\section{The \code{@render} Decorator}

Tools opt into generative UI with the \code{@render} decorator. It takes a component name from the built-in catalog:

\begin{lstlisting}
@render("component-name")
tool tool_name(params) -> ReturnType {
    """Description"""
    // body
}
\end{lstlisting}

The decorator does three things:

\begin{enumerate}
    \item Tells the compiler to include component metadata in the tool definition
    \item Tells the runtime to emit a \code{tool\_render} SSE event with the structured result
    \item Tells the chat UI to render the named component instead of showing raw text
\end{enumerate}

Tools without \code{@render} behave as before: the agent receives the result, interprets it, and streams a text response. Tools with \code{@render} still return their result to the agent for reasoning, but the UI additionally renders the component inline.

\begin{note}
The \code{@render} decorator is optional. Existing tools continue to work unchanged. Generative UI is additive.
\end{note}

\section{Built-in Component Catalog}

\haira{} ships a set of built-in components that cover common patterns. Each component expects a specific data shape and renders accordingly.

\subsection{table --- Tabular Data}

Renders a table with headers and rows. Supports optional column alignment and row highlighting.

\begin{lstlisting}
struct TableData {
    title: string
    headers: [string]
    rows: [[string]]
    highlight: [int]        // row indices to highlight (optional)
}

@render("table")
tool list_tables(file: string) -> TableData {
    """List all tables found in the Excel config file"""

    // ... extract tables ...
    return TableData{
        title: "Tables in " + filename,
        headers: ["Table", "Rows", "Type"],
        rows: [
            ["parameters", "42", "configuration"],
            ["batches", "8", "run"],
        ],
        highlight: [],
    }
}
\end{lstlisting}

\subsection{status-card --- Result Summary}

Renders a card with a status indicator (success, error, warning), title, and detail sections.

\begin{lstlisting}
struct StatusCard {
    status: string          // "success", "error", "warning", "info"
    title: string
    message: string
    sections: [CardSection]
}

struct CardSection {
    label: string
    content: string         // markdown supported
    style: string           // "default", "code", "error", "success"
}

@render("status-card")
tool validate(file: string) -> StatusCard {
    """Validate Excel data against the database schema"""

    // ... run validation ...
    return StatusCard{
        status: "error",
        title: "Validation Failed",
        message: "3 errors, 1 warning",
        sections: [
            CardSection{
                label: "Missing Columns",
                content: "`id_plant` in table `parameters`",
                style: "error",
            },
        ],
    }
}
\end{lstlisting}

\subsection{code-block --- Source Code or SQL}

Renders a syntax-highlighted code block with a copy button and optional title.

\begin{lstlisting}
struct CodeBlock {
    title: string           // optional header
    language: string        // "sql", "json", "haira", etc.
    code: string
}

@render("code-block")
tool generate_sql(file: string) -> CodeBlock {
    """Generate SQL INSERT statements from the config file"""

    // ... generate SQL ...
    return CodeBlock{
        title: "Generated Seeds SQL",
        language: "sql",
        code: sql_output,
    }
}
\end{lstlisting}

\subsection{diff --- Before/After Comparison}

Renders a side-by-side or unified diff view.

\begin{lstlisting}
struct DiffView {
    title: string
    before_label: string
    after_label: string
    before: string
    after: string
    language: string        // syntax highlighting language
}

@render("diff")
tool compare_schemas(table: string) -> DiffView {
    """Compare Excel schema with database schema for a table"""

    return DiffView{
        title: "Schema Diff: " + table,
        before_label: "Database",
        after_label: "Excel",
        before: db_schema,
        after: excel_schema,
        language: "sql",
    }
}
\end{lstlisting}

\subsection{key-value --- Property List}

Renders a compact list of labeled values. Useful for metadata, summaries, and configuration display.

\begin{lstlisting}
struct KeyValueList {
    title: string
    items: [KVItem]
}

struct KVItem {
    key: string
    value: string
    style: string           // "default", "success", "error", "muted"
}

@render("key-value")
tool file_info(file: string) -> KeyValueList {
    """Show metadata about the uploaded config file"""

    return KeyValueList{
        title: "File Info",
        items: [
            KVItem{ key: "Filename", value: filename, style: "default" },
            KVItem{ key: "Tables", value: "12", style: "default" },
            KVItem{ key: "Rows", value: "347", style: "default" },
            KVItem{ key: "Status", value: "Valid", style: "success" },
        ],
    }
}
\end{lstlisting}

\subsection{progress --- Multi-Step Progress}

Renders a vertical progress tracker showing completed, active, and pending steps. Useful for tools that report pipeline state.

\begin{lstlisting}
struct ProgressView {
    title: string
    steps: [ProgressStep]
}

struct ProgressStep {
    name: string
    status: string          // "done", "active", "pending", "failed"
    detail: string          // optional detail text
}

@render("progress")
tool pipeline_status(mr_id: string) -> ProgressView {
    """Check the deployment pipeline status"""

    return ProgressView{
        title: "Deployment Pipeline",
        steps: [
            ProgressStep{ name: "Branch Created", status: "done", detail: "" },
            ProgressStep{ name: "CI/CD Running", status: "active", detail: "2m 15s" },
            ProgressStep{ name: "Merge", status: "pending", detail: "" },
            ProgressStep{ name: "Deploy to QUA", status: "pending", detail: "" },
        ],
    }
}
\end{lstlisting}

\subsection{form --- Interactive Input}

Renders a form that the user can fill and submit. The submission triggers a follow-up tool call or message. This is the foundation for human-in-the-loop approval flows.

\begin{lstlisting}
struct FormView {
    title: string
    fields: [FormField]
    submit_label: string
    submit_action: string   // tool name or message prefix
}

struct FormField {
    name: string
    label: string
    field_type: string      // "text", "select", "checkbox", "textarea"
    value: string           // default value
    options: [string]       // for "select" type
    required: bool
}

@render("form")
tool request_approval(branch: string, mr_url: string) -> FormView {
    """Request user approval before pushing to GitLab"""

    return FormView{
        title: "Confirm Push to GitLab",
        fields: [
            FormField{
                name: "confirm",
                label: "Push branch '" + branch + "' and create MR?",
                field_type: "select",
                value: "",
                options: ["Yes, push it", "No, cancel"],
                required: true,
            },
            FormField{
                name: "comment",
                label: "MR comment (optional)",
                field_type: "textarea",
                value: "",
                options: [],
                required: false,
            },
        ],
        submit_label: "Submit",
        submit_action: "handle_approval",
    }
}
\end{lstlisting}

When the user submits the form, the chat sends a message in the format: \code{"[Form: handle\_approval] confirm=Yes, push it; comment=Looks good"}, which the agent can interpret and act upon.

\subsection{Component Summary}

\begin{table}[h]
\centering
\begin{tabular}{lll}
\toprule
\textbf{Component} & \textbf{Use Case} & \textbf{Data Shape} \\
\midrule
\code{table} & Tabular results, lists & Headers + rows \\
\code{status-card} & Success/error summaries & Status + sections \\
\code{code-block} & SQL, JSON, code output & Language + code string \\
\code{diff} & Before/after comparisons & Two text blocks \\
\code{key-value} & Metadata, config display & Key-value pairs \\
\code{progress} & Pipeline/step tracking & Steps with status \\
\code{form} & Approvals, user input & Fields + submit action \\
\bottomrule
\end{tabular}
\caption{Built-in generative UI components}
\label{tab:genui-components}
\end{table}

\section{Streaming Protocol}

Generative UI extends the existing SSE protocol with a new event type: \code{tool\_render}.

\subsection{Current Protocol}

The existing protocol emits three event types:

\begin{lstlisting}[language={},numbers=none]
event: tool_start
data: {"tool": "validate", "args": "{\"file\": \"/tmp/config.xlsx\"}"}

event: tool_end
data: {"tool": "validate", "ok": true}

data: {"delta": "The validation passed with..."}
\end{lstlisting}

\subsection{Extended Protocol}

When a tool has an \code{@render} decorator, the runtime emits \code{tool\_render} between \code{tool\_start} and \code{tool\_end}:

\begin{lstlisting}[language={},numbers=none]
event: tool_start
data: {"tool": "validate", "args": "{\"file\": \"/tmp/config.xlsx\"}"}

event: tool_render
data: {
    "tool": "validate",
    "component": "status-card",
    "props": {
        "status": "error",
        "title": "Validation Failed",
        "message": "3 errors, 1 warning",
        "sections": [...]
    }
}

event: tool_end
data: {"tool": "validate", "ok": false}
\end{lstlisting}

The \code{tool\_render} event carries:
\begin{itemize}
    \item \code{tool}: The tool name (for matching with the active tool card)
    \item \code{component}: The component name from the catalog
    \item \code{props}: The tool's return value, JSON-serialized, used as component props
\end{itemize}

The chat UI receives \code{tool\_render}, looks up the component in its registry, and renders it inline---replacing or augmenting the tool card.

\subsection{Dual Consumption}

The tool result is consumed twice:

\begin{enumerate}
    \item \textbf{By the UI}: The \code{tool\_render} event renders the component for the user
    \item \textbf{By the agent}: The result is still appended to the message history, so the agent can reason about it and decide the next step
\end{enumerate}

This means the agent sees the structured data (to make decisions) and the user sees the rich component (for readability). Neither path is sacrificed.

\section{Agent Behavior with Generative UI}

Agents do not need any changes to use generative UI. The \code{@render} decorator is on the tool, not the agent. An agent with a mix of rendered and non-rendered tools works naturally:

\begin{lstlisting}
agent Maltimize {
    provider: azure_openai
    system: "You help deploy Excel configs to PostgreSQL."
    tools: [
        read_excel,           // no @render -- text result
        validate,             // @render("status-card")
        generate_sql,         // @render("code-block")
        dry_run,              // @render("status-card")
        push_to_gitlab,       // no @render -- text result
    ]
    memory: conversation(max_turns: 20)
    temperature: 0.3
}
\end{lstlisting}

When the agent calls \code{validate}, the user sees a status card. When it calls \code{push\_to\_gitlab}, the user sees the agent's text interpretation. Both are valid---generative UI is opt-in per tool.

\section{Fallback Behavior}

Generative UI degrades gracefully in two scenarios:

\subsection{Non-Chat Contexts}

When a tool with \code{@render} is called outside a chat workflow (e.g., from a form workflow or a plain function), the decorator is ignored. The tool returns its structured data normally. The \code{@render} decorator only takes effect when the tool is executed through an agent's streaming loop.

\subsection{API Consumers}

When the SSE stream is consumed by a programmatic client (not the built-in chat UI), the \code{tool\_render} event provides the component name and props as JSON. The client can:
\begin{itemize}
    \item Render its own implementation of the component
    \item Ignore \code{tool\_render} events and use only \code{tool\_end} for status
    \item Use the \code{props} data directly for custom processing
\end{itemize}

\section{Complete Example}

A full migration assistant with generative UI:

\begin{lstlisting}
import "io"
import "http"
import "excel"
import "postgres"

provider azure_openai {
    type: "azure-openai"
    model: env("AZURE_OPENAI_DEPLOYMENT_NAME")
}

// --- Rendered tools ---

@render("table")
tool read_config(file_path: string) -> TableData {
    """Read an Excel config file and list the tables found"""

    wb = excel.open(file_path)?
    index = wb.read_sheet("index")?
    rows = []
    for entry in index {
        table = "" + entry["Table"]
        sheet = wb.read_sheet(get_sheet(table)) orelse []
        rows = rows + [[table, conv.int_to_string(len(sheet)), get_type(table)]]
    }
    wb.close()
    return TableData{
        title: "Config Tables",
        headers: ["Table", "Rows", "Type"],
        rows: rows,
        highlight: [],
    }
}

@render("status-card")
tool dry_run(file_path: string) -> StatusCard {
    """Dry-run the generated SQL against the local database"""

    // ... generate and test SQL ...
    if result["ok"] == true {
        return StatusCard{
            status: "success",
            title: "Dry Run Passed",
            message: "All SQL is valid",
            sections: [],
        }
    }
    return StatusCard{
        status: "error",
        title: "Dry Run Failed",
        message: result["error"],
        sections: [
            CardSection{
                label: "Root Cause",
                content: analysis["root_cause"],
                style: "error",
            },
            CardSection{
                label: "Fix SQL",
                content: "```sql\n" + analysis["fix"] + "\n```",
                style: "code",
            },
        ],
    }
}

// --- Agent ---

agent ConfigAssistant {
    provider: azure_openai
    system: "You deploy Excel configs to PostgreSQL via GitLab."
    tools: [read_config, dry_run, push_to_gitlab]
    memory: conversation(max_turns: 20)
    temperature: 0.3
}

// --- Workflow ---

@webui(title: "Config Assistant", description: "Chat-based config migration")
@webhook("/api/chat")
workflow Chat(message: string, session_id: string, file_path: file) -> stream {
    msg = message
    if file_path != "" { msg = "${message}\n\n[Attached file: ${file_path}]" }
    return ConfigAssistant.stream(msg, session: session_id)
}

fn main() {
    server = http.Server([Chat])
    server.listen(9001)
}
\end{lstlisting}

The user uploads an Excel file and says ``deploy this''. The agent calls \code{read\_config}---the chat renders a table showing all tables with row counts. Then calls \code{dry\_run}---the chat renders a status card showing success or failure with detailed sections. If it passes, the agent asks for confirmation before pushing.

All of this happens in a standard \haira{} program. No React, no Next.js, no bundler, no client-side framework.

\section{Custom Components}

The built-in catalog covers common patterns, but domain-specific use cases need custom components. \haira{} supports two levels of customization: composite components built from built-in primitives, and fully custom components written as TypeScript Web Components.

\subsection{Composite Components}

The simplest way to extend the UI is to compose built-in components. A composite component is a Haira struct whose fields map to multiple built-in components, rendered together as a group:

\begin{lstlisting}
@render("composite")
tool migration_report(file: string) -> MigrationReport {
    """Generate a full migration report with tables, validation, and SQL"""

    return MigrationReport{
        components: [
            Component{
                type: "key-value",
                props: KeyValueList{
                    title: "File Summary",
                    items: [
                        KVItem{ key: "File", value: filename, style: "default" },
                        KVItem{ key: "Tables", value: "12", style: "default" },
                    ],
                },
            },
            Component{
                type: "table",
                props: TableData{
                    title: "Tables",
                    headers: ["Name", "Rows", "Type"],
                    rows: table_rows,
                    highlight: [],
                },
            },
            Component{
                type: "status-card",
                props: StatusCard{
                    status: "success",
                    title: "Validation Passed",
                    message: "All checks passed",
                    sections: [],
                },
            },
        ],
    }
}
\end{lstlisting}

The \code{composite} renderer lays out multiple built-in components vertically in a single tool render. This requires no custom JavaScript---only data.

\subsection{Custom Web Components}

For fully custom rendering, users write a TypeScript Web Component in a \code{components/} directory next to their \haira{} source:

\begin{lstlisting}[language={}]
project/
  main.haira
  components/
    gantt-chart.ts
    schema-diagram.ts
\end{lstlisting}

A custom component file follows a simple contract:

\begin{lstlisting}[language={}]
// components/gantt-chart.ts

interface GanttProps {
    title: string;
    tasks: Array<{
        name: string;
        start: string;
        end: string;
        progress: number;
    }>;
}

export class HairaGanttChart extends HTMLElement {
    connectedCallback() {
        this.attachShadow({ mode: "open" });
    }

    // Called by the runtime with the tool's return value
    setProps(props: GanttProps) {
        const shadow = this.shadowRoot!;
        shadow.innerHTML = `
            <style>
                :host { display: block; }
                .chart {
                    padding: 1rem;
                    background: var(--haira-bg-card);
                    border: 1px solid var(--haira-border);
                    border-radius: var(--haira-radius);
                }
                .title {
                    font-weight: 600;
                    font-size: 0.85rem;
                    color: var(--haira-text);
                    margin-bottom: 0.75rem;
                }
                .bar-row {
                    display: flex;
                    align-items: center;
                    gap: 0.5rem;
                    margin: 0.35rem 0;
                }
                .bar-label {
                    font-size: 0.78rem;
                    color: var(--haira-text-dim);
                    min-width: 120px;
                }
                .bar-track {
                    flex: 1;
                    height: 20px;
                    background: var(--haira-bg-elevated);
                    border-radius: var(--haira-radius-sm);
                    overflow: hidden;
                }
                .bar-fill {
                    height: 100%;
                    background: var(--haira-gold);
                    border-radius: var(--haira-radius-sm);
                    transition: width 0.3s;
                }
                .bar-pct {
                    font-size: 0.7rem;
                    color: var(--haira-muted);
                    font-family: var(--haira-mono);
                    min-width: 35px;
                }
            </style>
            <div class="chart">
                <div class="title">${this.esc(props.title)}</div>
                ${props.tasks.map(t => `
                    <div class="bar-row">
                        <span class="bar-label">${this.esc(t.name)}</span>
                        <div class="bar-track">
                            <div class="bar-fill"
                                 style="width: ${t.progress}%"></div>
                        </div>
                        <span class="bar-pct">${t.progress}%</span>
                    </div>
                `).join("")}
            </div>
        `;
    }

    private esc(s: string): string {
        return s.replace(/&/g, "&amp;")
                .replace(/</g, "&lt;")
                .replace(/>/g, "&gt;");
    }
}

// Export the component class and its custom element tag name
export default {
    tag: "haira-gantt-chart",
    component: HairaGanttChart,
};
\end{lstlisting}

The component must:
\begin{enumerate}
    \item Extend \code{HTMLElement}
    \item Use Shadow DOM for style isolation
    \item Implement a \code{setProps(props)} method that receives the tool's return value as a typed object
    \item Export a default object with \code{tag} (the custom element name) and \code{component} (the class)
\end{enumerate}

The component is referenced in \code{@render} by its file name (without extension):

\begin{lstlisting}
struct GanttData {
    title: string
    tasks: [Task]
}

struct Task {
    name: string
    start: string
    end: string
    progress: int
}

@render("gantt-chart")
tool project_timeline(project_id: string) -> GanttData {
    """Show the project timeline as a Gantt chart"""

    return GanttData{
        title: "Q1 Roadmap",
        tasks: [
            Task{ name: "Backend API", start: "2025-01-01",
                  end: "2025-02-15", progress: 80 },
            Task{ name: "Frontend", start: "2025-02-01",
                  end: "2025-03-15", progress: 40 },
        ],
    }
}
\end{lstlisting}

\subsection{Component Resolution}

When the compiler encounters \code{@render("name")}, it resolves the component in this order:

\begin{enumerate}
    \item \textbf{Built-in catalog}: Check if \code{name} matches a built-in component (\code{table}, \code{status-card}, \code{composite}, etc.)
    \item \textbf{Local components}: Check if \code{components/name.ts} exists relative to the \code{.haira} source file
    \item \textbf{Compile error}: If neither matches, emit a diagnostic with the list of available component names
\end{enumerate}

Example diagnostic:

\begin{lstlisting}[language=bash,numbers=none]
error: unknown render component "gantt-chart"
  --> main.haira:42:9
   |
42 | @render("gantt-chart")
   |         ^^^^^^^^^^^^
   |
   = available: table, status-card, code-block, diff,
                 key-value, progress, form, composite
   = hint: create components/gantt-chart.ts to register
           a custom component
\end{lstlisting}

\subsection{Build Process}

The \haira{} compiler handles custom components as part of its normal build pipeline. The existing pipeline is:

\begin{lstlisting}[language=bash,numbers=none]
.haira source --> Lexer --> Parser --> Checker --> Codegen --> main.go
                                                                |
go-runtime/haira/ (linked via go.mod replace) ---+              |
                                                  |              v
                                          go build (produces binary)
\end{lstlisting}

The runtime's built-in UI is a pre-compiled JavaScript bundle (\code{ui/dist/haira-ui.js}) embedded in the Go binary via \code{//go:embed}. HTML templates contain a \code{<script>} tag that loads this bundle at \code{/\_ui/assets/haira-ui.js}.

Custom components extend this pipeline with one additional step:

\begin{lstlisting}[language=bash,numbers=none]
haira build main.haira

1. Parse, check, codegen as normal
2. Scan components/ directory for .ts files
3. Generate components/_register.ts (entry point):
      import c0 from "./gantt-chart.ts";
      import c1 from "./schema-diagram.ts";
      customElements.define(c0.tag, c0.component);
      customElements.define(c1.tag, c1.component);
4. Bundle components/_register.ts --> components.js
   (using bun build --minify --target browser)
5. Embed components.js as a Go string constant in
   the generated main.go:
      const customComponentsJS = `...minified JS...`
6. Register a second asset route:
      GET /_ui/assets/components.js
7. HTML templates include a second <script> tag
   after the built-in UI bundle
\end{lstlisting}

The generated \code{main.go} includes the custom bundle:

\begin{lstlisting}[language={}]
// Generated by haira build
package main

import "haira-go-runtime/haira"

// Custom components bundle (empty if no components/ dir)
const customComponentsJS = `(function(){class HairaGanttChart
extends HTMLElement{...}customElements.define("haira-gantt-
chart",HairaGanttChart);...})()`

func main() {
    server := haira.NewServer(workflows)
    server.SetCustomComponentsJS(customComponentsJS)
    server.Listen(9001)
}
\end{lstlisting}

The server serves the custom bundle at \code{/\_ui/assets/components.js} and the HTML template includes both scripts:

\begin{lstlisting}[language={}]
<script src="/_ui/assets/haira-ui.js"></script>
<script src="/_ui/assets/components.js"></script>
\end{lstlisting}

If no \code{components/} directory exists, the custom bundle is empty and the second \code{<script>} tag is omitted. Existing projects are unaffected.

\begin{note}
The compiler requires \code{bun} to be installed for TypeScript compilation of custom components. If \code{bun} is not found and custom components exist, the compiler emits a clear error with installation instructions. Projects without custom components do not require \code{bun}.
\end{note}

\subsection{Referencing Components from Tools}

The connection between a tool and its component is the \code{@render} decorator. The decorator name must match either a built-in component or the filename (without \code{.ts}) of a file in \code{components/}:

\begin{lstlisting}
// Built-in: name matches catalog entry
@render("table")
tool list_items(...) -> TableData { ... }

// Custom: name matches components/gantt-chart.ts
@render("gantt-chart")
tool project_timeline(...) -> GanttData { ... }

// Custom: name matches components/schema-diagram.ts
@render("schema-diagram")
tool show_schema(...) -> SchemaView { ... }
\end{lstlisting}

The compiler validates this mapping at compile time. The runtime connects them at render time: when a \code{tool\_render} SSE event arrives with \code{"component": "gantt-chart"}, the chat UI looks up the custom element registered under that name, creates an instance, and calls \code{setProps()} with the tool's result data.

The custom element tag name is defined by the component's \code{export default \{ tag: "..." \}} --- the convention is \code{haira-\{name\}} (e.g., \code{haira-gantt-chart} for \code{gantt-chart.ts}), but the developer controls it. The runtime matches by the \code{@render} name, not the tag name.

\begin{note}
Custom components run in the browser's Shadow DOM sandbox. They cannot access the parent page's DOM or interfere with other components. The \haira{} runtime provides theme CSS custom properties (\code{--haira-gold}, \code{--haira-bg}, etc.) that custom components can use for consistent styling.
\end{note}

\subsection{Theme Integration}

Custom components inherit \haira{}'s theme through CSS custom properties. The runtime injects these variables into the document root, making them available inside Shadow DOM:

\begin{lstlisting}[language={}]
/* Inside a custom component's Shadow DOM styles: */
:host {
    display: block;
    font-family: var(--haira-font);
    color: var(--haira-text);
}
.chart-bar {
    background: var(--haira-gold);
    border-radius: var(--haira-radius-sm);
}
.error { color: var(--haira-error); }
.success { color: var(--haira-success); }
\end{lstlisting}

Available CSS custom properties:

\begin{table}[h]
\centering
\begin{tabular}{ll}
\toprule
\textbf{Variable} & \textbf{Purpose} \\
\midrule
\code{--haira-bg} & Page background \\
\code{--haira-bg-card} & Card/panel background \\
\code{--haira-bg-elevated} & Elevated surface background \\
\code{--haira-border} & Border color \\
\code{--haira-gold} & Brand accent (gold) \\
\code{--haira-gold-light} & Lighter gold variant \\
\code{--haira-text} & Primary text color \\
\code{--haira-text-dim} & Secondary text color \\
\code{--haira-muted} & Muted/disabled text \\
\code{--haira-success} & Success state (green) \\
\code{--haira-error} & Error state (red) \\
\code{--haira-warn} & Warning state (yellow) \\
\code{--haira-info} & Info state (blue) \\
\code{--haira-font} & System font stack \\
\code{--haira-mono} & Monospace font stack \\
\code{--haira-radius} & Standard border radius \\
\code{--haira-radius-sm} & Small border radius \\
\bottomrule
\end{tabular}
\caption{CSS custom properties available to custom components}
\label{tab:genui-css-vars}
\end{table}

\subsection{Interactive Components}

Custom components can be interactive. When a user interacts with a component (clicks a button, selects a row, submits a form), the component dispatches a custom DOM event that the chat UI captures and sends as a follow-up message to the agent:

\begin{lstlisting}[language={}]
// Inside a custom component's setProps():
const btn = shadow.getElementById("approve-btn")!;
btn.addEventListener("click", () => {
    this.dispatchEvent(new CustomEvent("haira-action", {
        bubbles: true,
        composed: true,    // crosses Shadow DOM boundary
        detail: {
            action: "approve",
            data: { branch: "feature/config-123" },
        },
    }));
});
\end{lstlisting}

The chat UI listens for \code{haira-action} events and sends the action to the agent as a message: \code{"[Action: approve] \{\"branch\": \"feature/config-123\"\}"}. The agent can then interpret this and call the appropriate tool.

This pattern enables:
\begin{itemize}
    \item Clickable table rows that trigger drill-down queries
    \item Approval/reject buttons on deployment cards
    \item Filter controls that re-run a tool with different parameters
    \item Form submissions within custom components
\end{itemize}

\subsection{Example: Custom Schema Diagram}

A complete custom component for visualizing database table relationships:

\begin{lstlisting}[language={}]
// components/schema-diagram.ts

interface SchemaProps {
    title: string;
    tables: Array<{
        name: string;
        columns: Array<{
            name: string;
            type: string;
            pk: boolean;
            fk: boolean;
        }>;
    }>;
}

export class HairaSchemaDiagram extends HTMLElement {
    connectedCallback() {
        this.attachShadow({ mode: "open" });
    }

    setProps(props: SchemaProps) {
        const shadow = this.shadowRoot!;
        shadow.innerHTML = `
            <style>
                :host { display: block; }
                .diagram {
                    padding: 1rem;
                    background: var(--haira-bg-card);
                    border: 1px solid var(--haira-border);
                    border-radius: var(--haira-radius);
                    display: flex;
                    flex-wrap: wrap;
                    gap: 0.75rem;
                }
                .table-box {
                    border: 1px solid var(--haira-border);
                    border-radius: var(--haira-radius-sm);
                    min-width: 150px;
                }
                .table-name {
                    padding: 0.4rem 0.6rem;
                    background: var(--haira-gold);
                    color: #1a0e04;
                    font-weight: 700;
                    font-size: 0.8rem;
                    border-radius:
                        var(--haira-radius-sm)
                        var(--haira-radius-sm) 0 0;
                }
                .col {
                    padding: 0.25rem 0.6rem;
                    font-size: 0.75rem;
                    color: var(--haira-text-dim);
                    font-family: var(--haira-mono);
                    border-top: 1px solid var(--haira-border);
                }
                .col.pk { color: var(--haira-gold); }
                .col.fk { color: var(--haira-info); }
                .title {
                    font-weight: 600;
                    font-size: 0.85rem;
                    color: var(--haira-text);
                    width: 100%;
                    margin-bottom: 0.25rem;
                }
            </style>
            <div class="diagram">
                <div class="title">${this.esc(props.title)}</div>
                ${props.tables.map(t => `
                    <div class="table-box">
                        <div class="table-name">${this.esc(t.name)}</div>
                        ${t.columns.map(c => `
                            <div class="col ${c.pk ? "pk" : c.fk ? "fk" : ""}">
                                ${this.esc(c.name)}: ${this.esc(c.type)}
                            </div>
                        `).join("")}
                    </div>
                `).join("")}
            </div>
        `;
    }

    private esc(s: string): string {
        return s.replace(/&/g, "&amp;")
                .replace(/</g, "&lt;")
                .replace(/>/g, "&gt;");
    }
}

export default {
    tag: "haira-schema-diagram",
    component: HairaSchemaDiagram,
};
\end{lstlisting}

Used in \haira{}:

\begin{lstlisting}
struct SchemaView {
    title: string
    tables: [TableView]
}

struct TableView {
    name: string
    columns: [ColumnView]
}

struct ColumnView {
    name: string
    type: string
    pk: bool
    fk: bool
}

@render("schema-diagram")
tool show_schema(table_names: string) -> SchemaView {
    """Visualize the database schema for the given tables"""

    // ... query schema, build view ...
    return SchemaView{
        title: "Database Schema",
        tables: tables,
    }
}
\end{lstlisting}

The result: the agent calls \code{show\_schema}, and the chat renders an interactive diagram with color-coded primary and foreign keys---no React, no build system, just a \code{.ts} file next to the \code{.haira} source.

\section{Design Principles}

\begin{enumerate}
    \item \textbf{Tools own their rendering.} The \code{@render} decorator lives on the tool, not the agent or workflow. This keeps rendering co-located with the data source.

    \item \textbf{Structured data, not markup.} Tools return typed structs, not HTML or JSX. The runtime maps data to components. This keeps tools testable and platform-independent.

    \item \textbf{Catalog first, custom when needed.} The built-in catalog covers common patterns with zero configuration. Custom components extend the catalog with a single \code{.ts} file---no framework, no bundler, no build pipeline.

    \item \textbf{Additive, not breaking.} Every tool works without \code{@render}. Adding it improves the UI without changing the tool's logic or the agent's behavior.

    \item \textbf{Dual consumption.} The UI and the agent both receive the tool result. The UI renders a component; the agent reasons about the data. Neither path is secondary.

    \item \textbf{Progressive complexity.} Simple needs use built-in components (data only). Moderate needs use composites (multiple built-ins). Advanced needs use custom TypeScript Web Components (one \code{.ts} file). The learning curve matches the complexity of the use case.

    \item \textbf{Interactivity through events.} Custom components communicate back to the agent through a single event contract (\code{haira-action}). No callback registration, no bidirectional binding---just events that become messages.
\end{enumerate}

\section{External Frontend API}

The built-in chat UI is a convenience, not a requirement. \haira{} exposes a complete HTTP API that gives external frontends---React, Vue, Svelte, mobile apps, or any HTTP client---the same capabilities as the built-in UI. The built-in UI is itself just a consumer of this API.

\subsection{Discovery: \code{GET /\_api/}}

Returns metadata about all registered workflows. This is the same data the built-in UI uses to render the index page and route to form/chat views.

\begin{lstlisting}[language={},numbers=none]
GET /_api/

{
    "workflows": [
        {
            "name": "UploadConfig",
            "path": "/api/upload-config",
            "method": "POST",
            "ui_type": "form",
            "title": "Maltimize Config",
            "description": "Excel-to-PostgreSQL configuration migration",
            "params": [
                { "name": "file_path", "type": "file", "required": true },
                { "name": "debug", "type": "string", "required": false }
            ],
            "steps": [
                "Open Excel", "Extract table data", "Query schema",
                "Validate", "Generate SQL", "Dry-run SQL",
                "Push to GitLab", "Wait for CI/CD",
                "Wait for merge", "Verify deployment"
            ]
        },
        {
            "name": "Chat",
            "path": "/api/chat",
            "method": "POST",
            "ui_type": "chat",
            "title": "Maltimize Chat",
            "description": "Chat-based config migration assistant",
            "params": [
                { "name": "message", "type": "string", "required": true },
                { "name": "session_id", "type": "string", "required": true },
                { "name": "file_path", "type": "file", "required": false }
            ],
            "steps": [],
            "tools": [
                {
                    "name": "read_config_excel",
                    "description": "Read an Excel configuration file...",
                    "render": null
                },
                {
                    "name": "validate_config",
                    "description": "Validate Excel config data...",
                    "render": "status-card"
                },
                {
                    "name": "dry_run_migration",
                    "description": "Run the generated SQL...",
                    "render": "status-card"
                }
            ],
            "components": [
                {
                    "name": "status-card",
                    "type": "builtin"
                },
                {
                    "name": "gantt-chart",
                    "type": "custom",
                    "tag": "haira-gantt-chart"
                }
            ]
        }
    ]
}
\end{lstlisting}

The \code{tools} array lists each tool available to the agent, including its \code{render} component (or \code{null} if the tool has no UI). The \code{components} array lists all generative UI components used by the workflow, distinguishing built-in from custom.

\subsection{Workflow Metadata: \code{GET /\_api/\{path\}}}

Returns detailed metadata for a single workflow:

\begin{lstlisting}[language={},numbers=none]
GET /_api/api/chat

{
    "name": "Chat",
    "path": "/api/chat",
    "method": "POST",
    "ui_type": "chat",
    "title": "Maltimize Chat",
    "description": "Chat-based config migration assistant",
    "params": [...],
    "tools": [...],
    "components": [...]
}
\end{lstlisting}

\subsection{Streaming: \code{POST /\{path\}} with SSE}

The same endpoint used by the built-in UI. An external client sends a request with \code{Accept: text/event-stream} to receive SSE events:

\begin{lstlisting}[language=bash,numbers=none]
$ curl -N -H "Accept: text/event-stream" \
    -d '{"message":"deploy this","session_id":"abc123"}' \
    http://localhost:9001/api/chat
\end{lstlisting}

The SSE stream emits these event types:

\begin{table}[h]
\centering
\begin{tabular}{llp{7cm}}
\toprule
\textbf{Event} & \textbf{When} & \textbf{Data} \\
\midrule
\code{tool\_start} & Tool begins & \code{\{"tool": "name", "args": "\{...\}"\}} \\
\code{tool\_render} & Tool has UI result & \code{\{"tool": "name", "component": "status-card", "props": \{...\}\}} \\
\code{tool\_end} & Tool finishes & \code{\{"tool": "name", "ok": true\}} \\
(default) & Text delta & \code{\{"delta": "text chunk"\}} \\
\code{step} & Pipeline step event & \code{\{"name": "...", "status": "start|end|failed"\}} \\
\code{[DONE]} & Stream complete & --- \\
\bottomrule
\end{tabular}
\caption{SSE event types}
\label{tab:genui-sse-events}
\end{table}

An external frontend handles these events the same way the built-in UI does:

\begin{enumerate}
    \item On \code{tool\_start}: show a loading indicator for the tool
    \item On \code{tool\_render}: render the component using the \code{component} name and \code{props} data. For built-in components, the frontend implements its own version (e.g., a React \code{<StatusCard>}). For custom components, it uses the \code{props} JSON directly.
    \item On \code{tool\_end}: mark the tool as complete (success or failure)
    \item On text deltas: append to the assistant message
    \item On \code{[DONE]}: finalize the response
\end{enumerate}

\subsection{Form Submission: \code{POST /\{path\}} with JSON}

For non-streaming workflows (form UI), the same endpoint accepts JSON and returns JSON:

\begin{lstlisting}[language=bash,numbers=none]
$ curl -X POST http://localhost:9001/api/upload-config \
    -F "file_path=@config.xlsx" \
    -F "debug=true"

# Response with step-by-step SSE:
event: step
data: {"name": "Open Excel", "status": "start"}

event: step
data: {"name": "Open Excel", "status": "end", "duration_ms": 120}

# ... more steps ...

event: result
data: {"status": "success", "message": "Configuration deployed"}
\end{lstlisting}

Without \code{Accept: text/event-stream}, the response is plain JSON:

\begin{lstlisting}[language=bash,numbers=none]
$ curl -X POST http://localhost:9001/api/upload-config \
    -F "file_path=@config.xlsx"

{"status": "success", "message": "Configuration deployed to QUA"}
\end{lstlisting}

\subsection{CORS Support}

When an external frontend runs on a different origin (e.g., \code{localhost:3000} for a React dev server), the \haira{} server enables CORS automatically when the \code{HAIRA\_CORS\_ORIGIN} environment variable is set:

\begin{lstlisting}[language=bash,numbers=none]
$ HAIRA_CORS_ORIGIN="http://localhost:3000" ./maltimize
\end{lstlisting}

The wildcard \code{*} is also supported for development:

\begin{lstlisting}[language=bash,numbers=none]
$ HAIRA_CORS_ORIGIN="*" ./maltimize
\end{lstlisting}

\subsection{External Frontend Example}

A React frontend consuming the \haira{} API:

\begin{lstlisting}[language={}]
// React component consuming Haira's SSE API

function Chat() {
    const [messages, setMessages] = useState([]);
    const [tools, setTools] = useState(new Map());

    async function send(text: string, file?: File) {
        // Build request
        const body = new FormData();
        body.append("message", text);
        body.append("session_id", sessionId);
        if (file) body.append("file_path", file);

        // Open SSE connection
        const response = await fetch("/api/chat", {
            method: "POST",
            body,
            headers: { Accept: "text/event-stream" },
        });

        const reader = response.body.getReader();
        const decoder = new TextDecoder();

        while (true) {
            const { done, value } = await reader.read();
            if (done) break;

            const lines = decoder.decode(value).split("\n");
            for (const line of lines) {
                if (line.startsWith("event: tool_start")) {
                    // Show loading state for the tool
                }
                if (line.startsWith("event: tool_render")) {
                    const data = JSON.parse(/* next data line */);
                    // data.component = "status-card"
                    // data.props = { status: "error", ... }
                    // Render <StatusCard {...data.props} />
                }
                if (line.startsWith("data: ")) {
                    const data = JSON.parse(line.slice(6));
                    if (data.delta) {
                        // Append text to assistant message
                    }
                }
            }
        }
    }

    return (
        <div>
            {messages.map(msg =>
                msg.component
                    ? <ComponentRenderer
                          name={msg.component}
                          props={msg.props} />
                    : <MessageBubble text={msg.text} />
            )}
        </div>
    );
}

// Map Haira component names to React components
function ComponentRenderer({ name, props }) {
    const components = {
        "status-card": StatusCard,
        "table": DataTable,
        "code-block": CodeBlock,
        "gantt-chart": GanttChart,  // custom component
    };
    const Component = components[name];
    return Component ? <Component {...props} /> : null;
}
\end{lstlisting}

The external frontend renders its own components for each \code{tool\_render} event. The \code{component} name and \code{props} JSON provide all the information needed---the frontend maps component names to its own implementations (React, Vue, Svelte, or native mobile components).

\subsection{API Parity Guarantee}

The built-in UI has no privileged access. Every feature available in the built-in chat and form UIs is available through the HTTP API:

\begin{table}[h]
\centering
\begin{tabular}{lll}
\toprule
\textbf{Feature} & \textbf{Built-in UI} & \textbf{External API} \\
\midrule
Workflow discovery & Index page & \code{GET /\_api/} \\
Workflow metadata & \code{<script id="haira-meta">} & \code{GET /\_api/\{path\}} \\
Form submission & \code{<haira-form>} submit & \code{POST /\{path\}} (JSON) \\
Chat streaming & \code{<haira-chat>} SSE & \code{POST /\{path\}} (SSE) \\
Tool execution status & \code{<haira-tool-card>} & \code{tool\_start}/\code{tool\_end} events \\
Generative UI render & \code{<haira-*>} components & \code{tool\_render} event + props \\
File upload & Drag \& drop / attach & \code{multipart/form-data} \\
Step pipeline & \code{<haira-pipeline>} & \code{step} events \\
\bottomrule
\end{tabular}
\caption{API parity between built-in UI and external frontends}
\label{tab:genui-api-parity}
\end{table}

\begin{note}
The built-in UI can be disabled entirely with the \code{HAIRA\_DISABLE\_UI} environment variable. The API endpoints remain active. This is useful when deploying \haira{} as a pure backend behind a separate frontend application.
\end{note}

\section{Grammar Extension}

The \code{@render} decorator uses the existing decorator syntax. No new grammar is required:

\begin{lstlisting}[language=EBNF]
tool_decl     = { decorator } "tool" identifier "(" [ params ] ")" [ "->" type ]
                "{" triple_string { statement } "}" ;

decorator     = "@" identifier [ "(" decorator_args ")" ] ;
\end{lstlisting}

The compiler validates that the component name in \code{@render("name")} matches a known component from the catalog. Unknown component names produce a compile-time error.
